\chapter{Technology Mapping}
\label{ch:map}

Technology mapping turns a technology-independent AIG into a network of
implementable cells: $K$-input lookup tables (LUTs) for FPGAs, or gates drawn
from a standard-cell library for ASICs. ABC's mappers all share the same
skeleton --- enumerate small \emph{cuts} at every node, evaluate each cut's cost
under the target technology, and select one cut per node by dynamic programming
over several passes that trade delay against area --- but they differ in the
target and in the cost model. The flagship is the priority-cut FPGA mapper in
\pkg{src/map/if}, driving the commands \cmd{if} (logic-network flow) and
\cmd{\&if} (GIA flow).

\glance{The \cmd{if} mapper keeps at most \id{nCutsMax} priority cuts per node
(default \textbf{8}, switch \cmd{-C}), each cut a K-feasible set of leaves with
K = LUT size (switch \cmd{-K}). One node's cuts are formed by merging the cut
sets of its two fanins. Cuts are ranked lexicographically by a mode-dependent
key (delay-first vs.\ area-first); the winner becomes \id{CutBest}. Mapping runs
a fixed schedule of rounds: delay $\to$ area-flow $\to$ exact-area recovery.}

\section{The priority-cut FPGA mapper: \pkg{src/map/if}}

\subsection{Data model}

Four structures in \file{src/map/if/if.h} carry the mapper. \id{If_Man_t} is the
manager (objects, memory managers, truth-table tables, timing). \id{If_Obj_t} is
a node; besides fanins (\id{pFanin0}/\id{pFanin1}), level, and reference counts
it embeds its own best cut \id{CutBest} and a pointer \id{pCutSet} to the working
cut set. \id{If_Par_t} holds every user parameter. A \id{If_Set_t} is a set of
cut pointers (\id{ppCuts}) with a running count \id{nCuts} bounded by
\id{nCutsMax}.

The priority cut itself is compact --- a flexible array of leaf IDs preceded by
its cost fields:

\begin{lstlisting}[style=abccode]
struct If_Cut_t_
{
    float     Area;        // area (or area-flow) of the cut
    float     Edge;        // the edge flow
    float     Power;       // the power flow
    float     Delay;       // delay of the cut
    word      Config;      // configuration string
    int       iCutFunc;    // TT ID of the cut
    ...
    unsigned  Cost    : 12; // user's cost (<= IF_COST_MAX)
    unsigned  fCompl  :  1; // complemented attribute
    unsigned  fUser   :  1; // using the user's area and delay
    unsigned  fUseless:  1; // cannot be used in the mapping
    unsigned  fAndCut :  1; // matched with AND gate
    unsigned  nLimit  :  8; // the maximum number of leaves (= LUT size)
    unsigned  nLeaves :  8; // the number of leaves
    unsigned  decDelay: 16; // pin-to-pin decomposition delay
    int       pLeaves[0];
};
\end{lstlisting}

\id{nLimit} is set to \id{pPars->nLutSize} by \id{If_CutSetup}, so it \emph{is}
the LUT size K. \id{uSign} (a 32-bit one-hot-per-leaf-ID hash) enables fast
K-feasibility and containment pre-tests.

\subsection{Defaults from \id{If_ManSetDefaultPars}}

\begin{table}[htbp]
\centering
\small
\begin{tabularx}{\linewidth}{@{}llX@{}}
\toprule
Field & Default & Meaning \\
\midrule
\id{nLutSize}    & $-1$ & LUT size K; if unset, taken from LUT library \id{LutMax} (\cmd{-K}) \\
\id{nCutsMax}    & $8$  & priority cuts kept per node (\cmd{-C}) \\
\id{nFlowIters}  & $1$  & area-flow recovery rounds (\cmd{-F}) \\
\id{nAreaIters}  & $2$  & exact-area recovery rounds (\cmd{-A}) \\
\id{DelayTarget} & $-1$ & required-time target (auto) \\
\id{Epsilon}     & $0.005$ & float comparison tolerance \\
\id{fPreprocess} & $1$  & run the delay preprocessing rounds \\
\id{fArea}       & $0$  & pure area mode off \\
\id{fExpRed}     & $1$  & expand/reduce best cuts between rounds \\
\id{fEdge}       & $1$  & use edge (fanin-count) flow in ranking \\
\bottomrule
\end{tabularx}
\caption{Selected defaults set by \id{If_ManSetDefaultPars} (\file{src/map/if/ifCore.c}). \id{nLutSize} has no built-in default: \cmd{if} requires either \cmd{-K} or a loaded LUT library.}
\end{table}

\subsection{The round schedule: \id{If_ManPerformMapping}}

\id{If_ManPerformMapping} sets up the CI cut sets and cut memory, derives a
reverse topological order, then calls \id{If_ManPerformMappingComb}, which runs
the fixed sequence of passes below. Each pass is one call to
\id{If_ManPerformMappingRound} with a \emph{Mode} (delay $=0$, area-flow $=1$,
exact-area $=2$).

\begin{table}[htbp]
\centering
\small
\begin{tabularx}{\linewidth}{@{}lllX@{}}
\toprule
Pass (label) & Mode & Guard & Purpose \\
\midrule
``Delay''   & 0 & \id{fPreprocess \&\& !fArea} & first delay-optimal pass (\id{fFirst}) \\
``Delay-2'' & 0 & \id{fFancy}=1 & second delay pass, alternate sort \\
``Area''    & 0 & \id{fArea}=1 & delay-mode pass under area sort \\
``Flow''    & 1 & $\times$ \id{nFlowIters} & area-flow recovery \\
``Area''    & 2 & $\times$ \id{nAreaIters} & exact (ref/deref) area recovery \\
\bottomrule
\end{tabularx}
\caption{Round schedule in \id{If_ManPerformMappingComb}. After each area pass, if \id{fExpRed} is set, \id{If_ManImproveMapping} (\file{src/map/if/ifReduce.c}) expands/reduces cuts. \id{FinalDelay}/\id{FinalArea} are read back from \id{RequiredGlo}/\id{AreaGlo}.}
\end{table}

\id{If_ManPerformMappingRound} chooses the sort mode for the pass --- area
(\id{SortMode}=1) whenever \emph{Mode} is nonzero or \id{fArea} is set, the
``fancy'' delay variant (\id{SortMode}=2) when \id{fFancy} is set, otherwise
delay (\id{SortMode}=0) --- then visits every AND node in topological order,
calling \id{If_ObjPerformMappingAnd}, and \id{If_ObjPerformMappingChoice} for
representative nodes of a choice class (\id{pObj->fRepr}). It finishes by calling
\id{If_ManComputeRequired} to propagate required times for the next pass.

\subsection{Per-node dynamic programming: \id{If_ObjPerformMappingAnd}}

For one AND node, the mapper first refreshes the fanout estimate \id{EstRefs}
(exactly \id{nRefs} in delay mode; the blend $(2\cdot\id{EstRefs}+\id{nRefs})/3$
in the area modes), and in an area mode dereferences the currently selected best
cut so that recomputed areas are measured against the rest of the mapping. It
then enumerates the cross-product of the two fanins' cut sets:

\begin{lstlisting}[style=abccode]
If_ObjForEachCut( pObj->pFanin0, pCut0, i )
If_ObjForEachCut( pObj->pFanin1, pCut1, k )
{
    pCut = pCutSet->ppCuts[pCutSet->nCuts];
    // K-feasibility pre-test using the leaf-ID signatures
    if ( If_WordCountOnes(pCut0->uSign | pCut1->uSign) > p->pPars->nLutSize )
        continue;
    ...
    if ( !If_CutMergeOrdered( p, pCut0, pCut1, pCut ) )   // ordered union of leaves
        continue;
    ...
    if ( !p->pPars->fSkipCutFilter && If_CutFilter( pCutSet, pCut, fSave0 ) )
        continue;                                         // dominated / dominating
    ...
    pCut->Delay = If_CutDelay( p, pObj, pCut );
    if ( Mode && pCut->Delay > pObj->Required + p->fEpsilon && pCutSet->nCuts > 0 )
        continue;                                         // violates required time
    pCut->Area = (Mode == 2)? If_CutAreaDerefed(p,pCut) : If_CutAreaFlow(p,pCut);
    If_CutSort( p, pCutSet, pCut );                       // insert by priority
}
\end{lstlisting}

The union of leaves is built by \id{If_CutMergeOrdered} (a sorted merge that
fails once more than \id{nLimit}$=$K distinct leaves appear); \id{If_CutMerge} is
the unordered variant used when truth-table permutations are tracked
(\id{fUseTtPerm}). \id{If_CutFilter} drops any new cut whose leaves contain, or
are contained by, an existing cut. Truth tables and any user cost callbacks
(\id{pFuncCost}, \id{pFuncCell}) run before ranking. After the loop the best cut
is \id{ppCuts[0]}, copied into \id{CutBest}; a trivial single-leaf cut is
appended, and in the area modes the chosen cut is re-referenced.

\subsection{Cut ranking: \id{If_ManSortCompare} / \id{If_CutSort}}

\id{If_CutSort} does an incremental insertion sort that keeps at most
\id{nCutsMax} cuts, always holding the best cut at index~0. The ordering key is
the lexicographic comparator \id{If_ManSortCompare}, whose field priority
depends on \id{SortMode}:

\begin{table}[htbp]
\centering
\small
\begin{tabularx}{\linewidth}{@{}llX@{}}
\toprule
\id{SortMode} & Meaning & Lexicographic key (first difference wins) \\
\midrule
$0$ & delay      & \id{Delay}, \id{nLeaves}, \id{Area}, \id{Edge}, \id{Power}, \id{fUseless} \\
$1$ & area (flow/exact) & \id{Area}, \id{Edge}, \id{Power}, \id{nLeaves}, \id{fUseless} \\
$2$ & delay (fancy/old) & \id{Delay}, \id{fUseless}, \id{Area}, \id{Edge}, \id{Power}, \id{nLeaves} \\
\bottomrule
\end{tabularx}
\caption{Cut comparison order in \id{If_ManSortCompare} (\file{src/map/if/ifCut.c}), regular (non-power) branch. All floats are compared with tolerance \id{fEpsilon}. In delay mode a faster cut always wins and ties break toward fewer leaves; in area mode the cheaper cut wins first.}
\end{table}

Area-flow (\id{If_CutAreaFlow}) charges the cut's own LUT area
(\id{If_CutLutArea}, from \id{pLutLib->pLutAreas[nLeaves]} or $1.0$ for a unit
library) plus, for each leaf, the leaf's best-cut area amortised by its
\id{EstRefs}. The exact-area passes instead use \id{If_CutAreaDerefed}, which
reference-counts the actual cone (\id{If_CutAreaRef}/\id{If_CutAreaDeref}) to get
true shared area.

\subsection{Commands \cmd{if} and \cmd{\&if}}

\cmd{if} (\id{Abc_CommandIf}, \file{src/base/abci/abc.c}) maps a logic network;
\cmd{\&if} (\id{Abc_CommandAbc9If}) maps the current GIA via
\id{Gia_ManPerformMapping} in \file{src/aig/gia/giaIf.c}, which fills an
\id{If_Par_t}, builds an \id{If_Man_t}, and calls \id{If_ManPerformMapping}. Both
require a LUT size: either \cmd{-K num} or a LUT library
(\id{If_LibLut_t}, loaded with \cmd{read\_lut}); giving \cmd{-K} disables the
library. \cmd{-C num} sets the priority-cut count.

\section{Structural choices sharpen mapping}

Because \id{If_Obj_t} carries an equivalence pointer \id{pEquiv} and
\id{If_ObjPerformMappingChoice} folds all class members' cuts into one node's cut
set, \cmd{if} can map over an AIG that carries \emph{structural choices}. Choices
are alternative, functionally equivalent subgraphs (typically produced by DCH,
\cmd{dch}; see Chapter~\ref{ch:proof}); mapping across them lets the DP pick the
structure that yields the best cut, usually improving both depth and area
without changing functionality. When choices are present ABC disables
expand/reduce (\id{fExpRed}=0) for correctness.

\section{Standard-cell mapping: \pkg{src/map/mapper} and \pkg{src/map/amap}}

For ASIC targets ABC offers two structural mappers over a genlib library. Both
consume the same \id{Mio_Library_t} but optimise differently.

\begin{table}[htbp]
\centering
\small
\begin{tabularx}{\linewidth}{@{}lX@{}}
\toprule
\cmd{map} & Delay-oriented mapper in \pkg{src/map/mapper} (manager
\id{Map_Man_t}, i.e.\ \id{Map_ManStruct_t_} in \file{mapperInt.h}). Uses a
precomputed \emph{supergate} library (\id{Map_SuperLib_t}) and 6-input cuts with
N-canonical matching, then recovers area (\id{fAreaRecovery}). Entry
\id{Abc_CommandMap}. \\
\addlinespace
\cmd{amap} & Area-oriented mapper in \pkg{src/map/amap} (manager
\id{Amap_Man_t}). Matches library gates directly against cut functions and
prioritises area; supports AND/XOR/MUX decomposition of the subject graph. Entry
\id{Abc_CommandAmap}. \\
\bottomrule
\end{tabularx}
\caption{The two classic standard-cell mappers. \cmd{map} leads with arrival-time (delay), \cmd{amap} leads with area.}
\end{table}

\section{Reading libraries: \pkg{src/map/mio} and \pkg{src/map/super}}

Genlib libraries are parsed by \pkg{src/map/mio} into a \id{Mio_Library_t}
(\id{Mio_LibraryStruct_t_} in \file{mioInt.h}): a linked list of
\id{Mio_Gate_t}, each with area \id{dArea}, a Boolean formula \id{pForm}, a truth
table, and a list of \id{Mio_Pin_t} pin timings. The reader is registered as
both \cmd{read\_genlib} and \cmd{read\_library} (\id{Mio_CommandReadGenlib},
\file{src/map/mio/mio.c}). \pkg{src/map/super} precomputes \emph{supergates} ---
small gate clusters treated as single library entries --- via \cmd{super} /
\cmd{super2}; the result is loaded by \cmd{read\_super}
(\id{Map_CommandReadLibrary}) for the \cmd{map} flow.

\section{Liberty and gate sizing: \pkg{src/map/scl}}

\pkg{src/map/scl} handles real Liberty (\texttt{.lib}) data. \cmd{read\_lib}
(calling \id{Abc_SclReadLiberty}) parses a Liberty file into an \id{SC_Lib}
(struct \id{SC_Lib_} in \file{sclLib.h}), a set of \id{SC_Cell} entries each
holding \id{SC_Pin} timing/capacitance tables; same-function cells are linked
into area-ordered rings (\id{pNext}/\id{pPrev}/\id{pRepr}) for sizing. \cmd{stime}
runs static timing on the mapped netlist. The physically-aware flows operate on
an already-mapped design: \cmd{buffer} inserts buffering,
\cmd{upsize}/\cmd{dnsize} (and \cmd{minsize}/\cmd{maxsize}) resize gates along
the critical path against the Liberty timing model.

\section{Other mappers}

\begin{table}[htbp]
\centering
\small
\begin{tabularx}{\linewidth}{@{}lX@{}}
\toprule
\pkg{src/map/mpm}  & A configurable priority-cut technology mapper (an alternative/experimental \cmd{if}-style engine). \\
\pkg{src/map/fpga} & The older LUT mapper that predates \pkg{if}; historically the \cmd{fpga} command (registration currently disabled in \file{abc.c}). \\
\pkg{src/map/cov}  & Mapping a network into SOPs/ESOPs (sum-of-products covering). \\
\pkg{src/map/emap} & A multi-output gate (standard-cell) mapper, command \cmd{emap}. \\
\bottomrule
\end{tabularx}
\caption{Secondary mapping packages under \pkg{src/map}.}
\end{table}

\section{Summary}

\begin{table}[htbp]
\centering
\small
\begin{tabularx}{\linewidth}{@{}lllXl@{}}
\toprule
Mapper & Command & Target & File & Key struct \\
\midrule
Priority-cut LUT & \cmd{if}, \cmd{\&if} & FPGA LUTs & \file{src/map/if}, \file{src/aig/gia/giaIf.c} & \id{If_Man_t} \\
Delay std-cell   & \cmd{map}   & std cells & \file{src/map/mapper} & \id{Map_Man_t} \\
Area std-cell    & \cmd{amap}  & std cells & \file{src/map/amap}   & \id{Amap_Man_t} \\
Genlib reader    & \cmd{read\_genlib}/\cmd{read\_library} & library & \file{src/map/mio} & \id{Mio_Library_t} \\
Supergates       & \cmd{super}, \cmd{read\_super} & library & \file{src/map/super} & \id{Map_SuperLib_t} \\
Liberty / sizing & \cmd{read\_lib}, \cmd{buffer}, \cmd{upsize}/\cmd{dnsize} & std cells & \file{src/map/scl} & \id{SC_Lib} \\
Configurable     & --- & LUT/cells & \file{src/map/mpm} & \id{Mpm_Man_t} \\
Multi-output     & \cmd{emap}  & std cells & \file{src/map/emap} & --- \\
SOP/ESOP cover   & ---         & SOP/ESOP  & \file{src/map/cov} & --- \\
\bottomrule
\end{tabularx}
\caption{The ABC mapping packages at a glance.}
\end{table}

\subsection{REPL: LUT mapping}

\begin{lstlisting}[style=abcshell]
abc 01> read i10.aig
abc 02> strash
abc 03> dch            # build structural choices (optional, improves mapping)
abc 04> if -K 6 -C 12  # 6-LUT mapping, keep 12 priority cuts per node
abc 05> print_stats    # reports LUT count and depth
\end{lstlisting}

\subsection{REPL: standard-cell mapping}

\begin{lstlisting}[style=abcshell]
abc 01> read_genlib mcnc.genlib   # or read_library
abc 02> read i10.aig
abc 03> strash
abc 04> map                       # delay-oriented standard-cell mapping
abc 05> print_gates               # gate usage / area report
\end{lstlisting}
